%! Vector-Valued Functions — Unit 4, Lesson 4.9 — Warm-Up (Answer Key)
\documentclass[10pt]{article}
\usepackage{precalculus-article}
\usepackage{precalculus-key}   % pulls in -boxes; do NOT also load -boxes
\newcommand{\TallMath}[1]{$\displaystyle #1\rule[-1.4em]{0pt}{3.2em}$}

\begin{document}

\pageheader{Unit 4, Lesson 4.9}{Warm-Up --- Answer Key}
\namedateperiod

\vspace{0.12in}

\begin{notesbox}{Spiral Review --- Key}
\small

\textbf{1.}\quad Given the parametric function $f(t)=(2t,\;t^{2}-1)$, find $f(3)$.

\vspace{0.08in}
\hspace{1em} $x(3) = $ \ans{$6$} \qquad $y(3) = $ \ans{$8$}
\qquad $f(3) = $ \ans{$(6,\,8)$}

\vspace{0.18in}

\textbf{2.}\quad Find the magnitude of the vector $\langle 3,\,-4\rangle$.

\vspace{0.08in}
\hspace{1em} $\|\langle 3,\,-4\rangle\| = \sqrt{3^{2}+(-4)^{2}} = \sqrt{9+16} = \sqrt{25} = $ \ans{$5$}

\begin{teachernote}
  Common error: students compute $\sqrt{3+(-4)}=\sqrt{-1}$ (forgetting to square).
  Remind: magnitude always uses \emph{squares} of the components.
\end{teachernote}

\vspace{0.18in}

\textbf{3.}\quad A particle is at position $(0,\,5)$ when $t=0$ and at position
$(6,\,5)$ when $t=3$.
Find the average rate of change of $x(t)$ over the interval $0\le t\le 3$.

\vspace{0.08in}
\hspace{1em} $\dfrac{\Delta x}{\Delta t} = \dfrac{6-0}{3-0} = \dfrac{6}{3} = $ \ans{$2$ units per unit time}

\end{notesbox}

\end{document}
